\documentclass[11pt,a4paper]{article}
\usepackage[T1]{fontenc}
\usepackage{lmodern}
\usepackage[margin=27mm]{geometry}
\usepackage{amsmath,amssymb,amsthm,mathtools}
\usepackage{microtype}
\usepackage{booktabs,longtable,array}
\usepackage{xcolor}
\usepackage[colorlinks=true,linkcolor=blue!45!black,citecolor=blue!45!black,
urlcolor=blue!45!black,breaklinks=true]{hyperref}
\usepackage{enumitem}
\usepackage{fancyhdr}
\hypersetup{pdftitle={An Ordered Poisson Coupling for Mittag-Leffler Laws},
pdfauthor={Research note prepared for Keio},
pdfsubject={Exact stochastic comparison, optimal scale, and quantitative CDF separation}}
\setlist{itemsep=4pt,topsep=6pt}
\setlength{\emergencystretch}{2em}
\setlength{\headheight}{14pt}
\pagestyle{fancy}
\fancyhf{}
\fancyhead[L]{\small Poisson coupling and Mittag--Leffler laws}
\fancyhead[R]{\small Research draft}
\fancyfoot[C]{\thepage}
\newtheorem{theorem}{Theorem}[section]
\newtheorem{proposition}[theorem]{Proposition}
\newtheorem{corollary}[theorem]{Corollary}
\DeclareMathOperator{\Var}{Var}
\title{\textbf{An Ordered Poisson Coupling\\for Mittag--Leffler Laws}\\[5pt]
\large Exact stochastic comparison, optimal scale, and a quantitative gap}
\author{Research note prepared for Keio}
\date{8 November 2026}
\begin{document}
\maketitle
\begin{abstract}
Let \(Z_\alpha\) have Laplace transform
\(\mathbb E e^{-tZ_\alpha}=e^{-t^\alpha}\), and set
\(M_\alpha=Z_\alpha^{-\alpha}\), \(0<\alpha<1\).
We give a complete probabilistic proof that
\(M_\alpha/\Gamma(1-\alpha)\) decreases in the usual stochastic order,
with optimal scale. This is the comparison proposed in Simon's 2014
Remark~9 and restated in Ferreira--Simon's 2025 Remark~7(b), with the
original normalization and full index range.
One Poisson arrival sequence realizes every member of the family as
\[
Y_\alpha=\left(\sum_{j\geq1}T_j^{-1/\alpha}\right)^{-\alpha}.
\]
The resulting trajectories are smooth, strictly decreasing, and
strictly log-concave. We also prove an explicit positive lower bound
for the CDF difference at every positive threshold, and a stronger
bound linear in the index difference on compact index intervals.
The Poisson-series and norm principles are classical and are credited
as such. Historical priority of the specific results remains unverified.
This proof draft is not presented as a certified success under all six
novelty requirements recorded in Appendix~A.
\end{abstract}
\section{Original statement and sources}
The stochastic-order clause was already conjectured by T. Simon,
\emph{Comparing Fr\'echet and positive stable laws}, Electronic Journal
of Probability 19 (2014), no. 16, 1--25, Theorem 3, equation (5.4),
and Remark 9; \cite{Simon}. The arXiv v2 is dated
27 January 2014. It conjectures that for all $0<\beta<\alpha<1$,
\[
 \Gamma(1-\beta)M_\alpha
     \prec_{\mathrm{st}}\Gamma(1-\alpha)M_\beta,
\]
where $X\prec_{\mathrm{st}}Y$ means $X\leq_{\mathrm{st}}Y$ and no
$c>1$ satisfies $cX\leq_{\mathrm{st}}Y$. The source proves a restricted
parameter region only. The unrestricted convex-order clause next to
it is a different assertion and is not counted as a new result here.
R. A. C. Ferreira and T. Simon, \emph{Mittag--Leffler functions and convex
ordering}, arXiv:2512.04940v1, 4 December 2025, Remark 7(b), propose that
\[
  \alpha\longmapsto \frac{M_\alpha}{\Gamma(1-\alpha)}
\]
is non-increasing for the usual stochastic order on $(0,1)$.
Here $M_\alpha=Z_\alpha^{-\alpha}$ in distribution, where
$\mathbb E e^{-tZ_\alpha}=e^{-t^\alpha}$ for $t\geq0$.
The current arXiv record, checked on 8 November 2026, lists only v1 \cite{FS}.
The author's current publications page also lists that preprint.
The original claim, normalization, index interval, and probability order
are retained exactly below.

\section{Main result and classical inputs}
The Poisson Laplace functional and monotonicity of $\ell^p$ norms are
classical. The series used below is the positive LePage series. None of
these ingredients is claimed to be new. The proposed contribution is
their simultaneous, correctly normalized application to the specific
stochastic-order conjecture, together with the explicit separation bound.
Poisson representations through \(\ell^p\) norms also occur explicitly in
Oesting's max-stable constructions \cite{Oesting}; that viewpoint is credited
as prior work. Jan\ss en's Corollary~6.3 \cite{Janssen}, reproducing
a result of Schlather, also identifies stable-power laws as limits of
sample norms. These are close precedents for the qualitative mechanism.
Standard Poisson references are \cite{LWZ,LP}. Priority of the
specific comparison or the quantitative estimates has not been established
by the bounded search described in Appendix~A.

\begin{theorem}[Full comparison and exact optimal scale]\label{thm:main}
There is a coupling \((Y_\alpha)_{0<\alpha<1}\) with
\[
Y_\alpha\stackrel d=\frac{M_\alpha}{\Gamma(1-\alpha)}
\]
and one event of probability one on which
\[
Y_\beta<Y_\alpha\qquad\hbox{for every }0<\alpha<\beta<1.
\]
For every such pair and every \(x>0\),
\[
\mathbb P(Y_\beta\leq x)>\mathbb P(Y_\alpha\leq x).
\]
Furthermore, for \(c>0\),
\[
cY_\beta\leq_{\mathrm{st}}Y_\alpha\quad\Longleftrightarrow\quad c\leq1.
\]
Equivalently, the largest \(c\) for which
\(cM_\beta\leq_{\mathrm{st}}M_\alpha\) is
\(\Gamma(1-\alpha)/\Gamma(1-\beta)\).
\end{theorem}
The following sections prove every clause. In particular, strictness
at every threshold is proved by an explicit positive-probability
event, not inferred from strictness of a coupling alone.

\section{Construction and marginal identification}\label{sec:construction}
Let $0<T_1<T_2<\cdots$ be the arrivals of one rate-one Poisson process,
and use this same sequence for every $0<\alpha<1$. On the probability-one
event $T_j/j\to1$, define
\[
 S_\alpha=\sum_{j\geq1}T_j^{-1/\alpha},\qquad
 Y_\alpha=S_\alpha^{-\alpha}.
\]
Every series converges on this one event, simultaneously for all
$\alpha\in(0,1)$: for each such $\alpha$, its tail is bounded by a
constant times $\sum j^{-1/\alpha}$.
For $t\geq0$, the Poisson Laplace functional gives
\begin{align*}
 \mathbb E e^{-tS_\alpha}
 &=\exp\left(-\int_0^\infty(1-e^{-t u^{-1/\alpha}})\,du\right)\\
 &=\exp\left(-\alpha t^\alpha
        \int_0^\infty(1-e^{-v})v^{-\alpha-1}\,dv\right)\\
 &=\exp\bigl(-\Gamma(1-\alpha)t^\alpha\bigr).
\end{align*}
The substitution is $v=t u^{-1/\alpha}$ and the last equality follows
by integration by parts; both boundary terms vanish for $0<\alpha<1$.
Thus
\[
 S_\alpha\stackrel d=\Gamma(1-\alpha)^{1/\alpha}Z_\alpha,
 \qquad
 Y_\alpha\stackrel d=\frac{M_\alpha}{\Gamma(1-\alpha)}.
\]
The Poisson Laplace functional here can be derived first for simple
nonnegative functions supported in a finite interval from independence
of the Poisson counts, then passed to this integrand by monotone
convergence. Consequently no distributional assertion about the
LePage series has to be imported separately.

\section{Pathwise order}
Write $a_j=T_j^{-1}$. Then
\[
 Y_\alpha=\|a\|_{1/\alpha}^{-1}.
\]
If $0<\alpha<\beta<1$, set $p=1/\alpha>q=1/\beta>1$.
Normalize $b_j=a_j/\|a\|_q$. Every $b_j$ lies in $(0,1)$ and
$\sum b_j^q=1$, so $\sum b_j^p<1$. Therefore
\[
       Y_\beta<Y_\alpha\quad\hbox{almost surely}.
\]
This proves the original usual stochastic order, with all index pairs
covered by a single coupling and one probability-one event.
Here \(X\leq_{\mathrm{st}}W\) means
\(\mathbb P(X>x)\leq\mathbb P(W>x)\) for every real \(x\).
The identity with a norm supplies the representation; the strict norm
inequality supplies the actual comparison.

\section{Optimality in Simon's original sense}
For $0<s<1$, the fractional-moment integral for a nonnegative random
variable yields
\[
 \mathbb E Z_\alpha^{\alpha s}
 =\frac{\alpha s}{\Gamma(1-\alpha s)}
   \int_0^\infty(1-e^{-t^\alpha})t^{-\alpha s-1}\,dt
 =\frac{\Gamma(1-s)}{\Gamma(1-\alpha s)}.
\]
The substitution $u=t^\alpha$ and integration by parts justify the
last step. Thus
\[
 \mathbb E Y_\alpha^{-s}
 =\frac{\Gamma(1-\alpha)^s\Gamma(1-s)}{\Gamma(1-\alpha s)},
 \qquad
 \lim_{s\uparrow1}
  \frac{\mathbb E Y_\beta^{-s}}{\mathbb E Y_\alpha^{-s}}=1.
\]
If $cY_\beta\leq_{\mathrm{st}}Y_\alpha$ held with $c>1$, applying
the decreasing function $y\mapsto y^{-s}$ would give
$c^{-s}\mathbb E Y_\beta^{-s}\geq\mathbb E Y_\alpha^{-s}$.
Letting $s\uparrow1$ gives $c^{-1}\geq1$, a contradiction.
Consequently $Y_\beta\prec_{\mathrm{st}}Y_\alpha$, which is exactly
Simon's original optimal order after positive rescaling.

\section{An explicit positive gap at every threshold}\label{sec:twopoint}
Fix $0<\alpha<\beta<1$ and $x>0$. Set
\[
 u=x2^\alpha,\qquad v=x2^\beta,\qquad
 a=\frac{2u+v}{3},\qquad b=\frac{u+2v}{3},\qquad
 A=x^{-1/\alpha}-2a^{-1/\alpha}>0.
\]
Choose, for example,
\[
 R=1+\max\left\{b,
      \left(\frac{2\alpha}{(1-\alpha)A}\right)^{\alpha/(1-\alpha)}
              \right\}.
\]
The expected tail of $S_\alpha$ beyond $R$ is
\[
 \mathbb E\sum_{T_j>R}T_j^{-1/\alpha}
       =\frac{\alpha}{1-\alpha}R^{1-1/\alpha}<A/2.
\]
By Markov's inequality the tail is less than $A$ with probability
greater than $1/2$. Independently, the event that there are exactly two
points in $[a,b]$ and no other points in $[0,R]$ has probability
$e^{-R}(b-a)^2/2$. On the intersection,
\[
 S_\alpha<2a^{-1/\alpha}+A=x^{-1/\alpha},\qquad
 S_\beta\geq2b^{-1/\beta}>x^{-1/\beta}.
\]
Hence $Y_\beta<x<Y_\alpha$ there, and the common coupling yields the
explicit bound
\[
 \mathbb P(Y_\alpha>x)-\mathbb P(Y_\beta>x)
       \geq\frac{e^{-R}(b-a)^2}{4}>0.
\]
This closes the stronger statement of strict stochastic separation for
every positive threshold, without relying on the generally invalid
inference that a strict coupling alone implies strict tail inequalities
at every threshold.

\section{Regularity, entropy, and curvature}\label{sec:entropy}
For a fixed sample on $T_j/j\to1$, the series and all its parameter
derivatives converge locally uniformly on $(0,1)$. Indeed, on a compact
index interval with upper endpoint $c<1$, they are dominated by a
constant times $j^{-1/c}(1+\log j)^k$ for each fixed derivative order.
Let
\[
 p_j(\alpha)=\frac{T_j^{-1/\alpha}}{S_\alpha},\qquad
 H(\alpha)=-\sum_jp_j(\alpha)\log p_j(\alpha).
\]
Direct differentiation gives
\[
 \frac{d}{d\alpha}\log Y_\alpha=-H(\alpha)<0,
 \qquad
 \frac{d^2}{d\alpha^2}\log Y_\alpha
   =-\frac{\operatorname{Var}_{p(\alpha)}(\log T_j)}{\alpha^3}<0.
\]
All sums are finite by the same domination. Strictness holds since
infinitely many weights are positive and the arrival times are distinct.
The coupling therefore has smooth, strictly decreasing, strictly
log-concave trajectories on $(0,1)$. This is a statement about sample
paths in the chosen coupling, not about log-concavity of the marginal
density in its spatial variable.

For an explicit finite-parameter identity, set \(r=\alpha/\beta\in(0,1)\).
Then
\[
\frac{Y_\beta}{Y_\alpha}
=\left(\sum_j p_j(\alpha)^r\right)^{-\beta}
=\exp\{-(\beta-\alpha)H_r(p(\alpha))\},
\qquad
H_r(p)=\frac{\log\sum_jp_j^r}{1-r}.
\]
Indeed,
\(S_\beta=S_\alpha^r\sum_jp_j(\alpha)^r\).
The sum is finite because both \(S_\alpha\) and \(S_\beta\) are finite.
Strict concavity also yields, for \(0<\theta<1\),
\[
Y_{\theta\alpha+(1-\theta)\beta}
>Y_\alpha^\theta Y_\beta^{1-\theta}.
\]

\section{Endpoints and Mittag--Leffler consequence}
The standard $\ell^p$ limit gives $Y_\alpha\to T_1$ as
$\alpha\downarrow0$. As $\alpha\uparrow1$, for every fixed $k$ the
norm is bounded below by
$(\sum_{j=1}^k T_j^{-1/\alpha})^\alpha\to\sum_{j=1}^kT_j^{-1}$.
Since $T_j/j\to1$, the latter partial sums diverge, proving
$Y_\alpha\to0$ almost surely.
Consequently, for $x>0$ and $0<\alpha<\beta<1$,
\[
 \frac1{1+x}
 < E_\alpha\!\left(-\frac{x}{\Gamma(1-\alpha)}\right)
 < E_\beta\!\left(-\frac{x}{\Gamma(1-\beta)}\right)<1.
\]
Here $E_\alpha(z)=\sum_{n\geq0}z^n/\Gamma(1+\alpha n)$.
The equality of its Laplace representation to that of $M_\alpha$ follows
from the negative moments of $Z_\alpha$:
\[
 \mathbb E Z_\alpha^{-r}
   =\frac1{\Gamma(r)}\int_0^\infty t^{r-1}e^{-t^\alpha}\,dt
   =\frac{\Gamma(1+r/\alpha)}{\Gamma(1+r)}\quad(r>0).
\]
At \(r=\alpha n\), these moments give
\(\mathbb E M_\alpha^n=n!/\Gamma(1+\alpha n)\).
For \(t\geq0\), Tonelli identifies
\(\mathbb E e^{tM_\alpha}\) with the defining series \(E_\alpha(t)\),
which is finite by Stirling's formula. Domination by
\(e^{|z|M_\alpha}\) then justifies the power-series calculation for
all complex \(z\), including the negative arguments used above.
The strict lower comparison uses
\(Y_\alpha<T_1\), since \(S_\alpha>T_1^{-1/\alpha}\).
The endpoint limits in the Laplace inequality follow by bounded
convergence. The value at $x=0$ is equality throughout.

\section{A locally uniform linear separation estimate}\label{sec:linear}
The preceding two-point event proves a positive gap for every pair.
The next result is stronger when the indices are close: its lower bound
is linear in their separation on each fixed compact index interval.

\begin{theorem}[Linear CDF separation]\label{thm:linear}
Fix $0<a<b<1$ and $x>0$. For every
$a\leq\alpha<\beta\leq b$, write
$F_\gamma(x)=\mathbb P(Y_\gamma\leq x)$. Define
\[
 C=1+\max\left\{2^b,
       \left(\frac{8xb}{1-b}\right)^{b/(1-b)}\right\},
 \qquad
 \eta=\left(\frac{x}{Cx+1}\right)^{1/a}.
\]
Then $0<\eta<1/2$ and
\[
 F_\beta(x)-F_\alpha(x)
 \geq (\beta-\alpha)\frac{x e^{-Cx}}4
             \log\frac1{1-\eta}>0.
\]
The constant is uniform over every pair in the chosen compact index
interval. In particular, unlike the preceding two-point-event estimate,
it remains linear as $\beta-\alpha$ decreases to zero.

\end{theorem}
\begin{proof}
Use independent rate-one exponential variables $E_1,E_2,\ldots$, and
put $V_1=0$, $V_j=E_2+\cdots+E_j$ for $j\geq2$.
The positive $V_j$ form a rate-one Poisson process independent of $E_1$.
Conditioning on all the $V_j$, define for $t>0$
\[
 S_\gamma(t)=\sum_{j\geq1}(t+V_j)^{-1/\gamma},
 \qquad Y_\gamma(t)=S_\gamma(t)^{-\gamma}.
\]
Almost surely $V_j/j\to1$. The sums and all parameter derivatives
converge locally uniformly on $(0,1)\times(0,\infty)$, by comparison
with $j^{-1/b'}(1+\log j)^k$ on any compact index range ending at
$b'<1$. The map $t\mapsto Y_\gamma(t)$ is smooth and strictly
increasing. It tends to zero as $t\downarrow0$, since
$Y_\gamma(t)\leq t$. It tends to infinity as $t\to\infty$, since
$S_\gamma(t)\to0$ by dominated convergence, using the summable
sequence $(1+V_j)^{-1/\gamma}$ for $t\geq1$.
Consequently a unique smooth function $\tau_\gamma(x)>0$ satisfies
$Y_\gamma(\tau_\gamma(x))=x$, and conditioning gives
\[
 F_\gamma(x)=1-\mathbb E e^{-\tau_\gamma(x)}.
\]
At $t=\tau_\gamma(x)$, let
\[
 p_j=\frac{(t+V_j)^{-1/\gamma}}{S_\gamma(t)},
 \qquad H=-\sum_jp_j\log p_j,
 \qquad D=\sum_j\frac{p_j}{t+V_j}.
\]
The previously proved entropy derivative and direct differentiation in
$t$ give $\partial_\gamma\log Y=-H$ and $\partial_t\log Y=D$.
The implicit-function theorem therefore gives
\[
 \frac{d}{d\gamma}\tau_\gamma(x)=\frac H D
       \geq\tau_\gamma(x)H,
\]
because $D\leq1/t$.

Consider the event
\[
 G=\left\{V_2\leq1,\quad
 \sum_{j\geq2}(C+V_j/x)^{-1/b}\leq1-C^{-1/b}\right\}.
\]
Campbell's formula gives
\[
 \mathbb E\sum_{j\geq2}(C+V_j/x)^{-1/b}
 =\frac{xb}{1-b}C^{1-1/b}<\frac18.
\]
The threshold $1-C^{-1/b}$ exceeds $1/2$. Markov's inequality
therefore bounds the failure probability of the second condition by
$1/4$. Since $\mathbb P(V_2\leq1)=1-e^{-1}>1/2$,
the union bound gives $\mathbb P(G)>1/4$; no independence between the
two conditions defining $G$ is assumed.

On $G$, uniformly for $\gamma\in[a,b]$, the fact that
$C+V_j/x>1$ implies
\[
 x^{1/\gamma}S_\gamma(Cx)
 =C^{-1/\gamma}+\sum_{j\geq2}(C+V_j/x)^{-1/\gamma}
 \leq C^{-1/b}+\sum_{j\geq2}(C+V_j/x)^{-1/b}\leq1.
\]
Thus $Y_\gamma(Cx)\geq x$. Together with $Y_\gamma(t)\leq t$,
this shows $x\leq\tau_\gamma(x)\leq Cx$. At the threshold root,
\[
 p_1=\left(\frac{x}{\tau_\gamma(x)}\right)^{1/\gamma}\geq\eta,
 \qquad
 p_2=\left(\frac{x}{\tau_\gamma(x)+V_2}\right)^{1/\gamma}\geq\eta.
\]
Every $p_j$ is therefore at most $1-\eta$, and hence
$H\geq\log(1/(1-\eta))$. It follows on $G$ that
\[
 e^{-\tau_\gamma(x)}\frac{d}{d\gamma}\tau_\gamma(x)
 \geq x e^{-Cx}\log\frac1{1-\eta},
\]
Integrate pathwise from $\alpha$ to $\beta$, restrict expectation to
$G$, and use nonnegativity on the complement. This proves the claimed
bound. No differentiation under an expectation is required.
\end{proof}

\section{A stronger endpoint normalization and a closed-form check}
\begin{proposition}\label{prop:endpoint}
In the same coupling, almost surely
\[
\lim_{\alpha\uparrow1}\frac{Y_\alpha}{1-\alpha}=1.
\]
Consequently \(\Gamma(1-\alpha)Y_\alpha\to1\) almost surely.
\end{proposition}
\begin{proof}
Put \(p=1/\alpha>1\). For any \(\varepsilon\in(0,1)\), the event
\(T_j/j\to1\) supplies \(J\) such that
\((1-\varepsilon)j\leq T_j\leq(1+\varepsilon)j\) for \(j\geq J\).
Integral comparison gives
\((p-1)\sum_{j\geq J}j^{-p}\to1\) as \(p\downarrow1\).
The finite prefix vanishes after multiplication by \(p-1\).
Squeezing and then letting \(\varepsilon\downarrow0\) proves
\[
(p-1)\sum_jT_j^{-p}\longrightarrow1.
\]
It follows that
\[
\log\frac{Y_\alpha}{1-\alpha}
=(\alpha-1)\log(1-\alpha)-\alpha\log\alpha+o(1)
\longrightarrow0.
\]
Finally \((1-\alpha)\Gamma(1-\alpha)\to1\).
\end{proof}

At \(\alpha=1/2\), the density of \(Y_{1/2}\) is
\[
f_{1/2}(y)=e^{-\pi y^2/4},\qquad y>0.
\]
Its integral is one and its mean is \(2/\pi\), consistent with the
moment formula. Completing the square in its Laplace integral yields
\[
\mathbb E e^{-xY_{1/2}}
=e^{x^2/\pi}\operatorname{erfc}(x/\sqrt{\pi})
=E_{1/2}(-x/\sqrt{\pi}).
\]
The last identity follows also by separating the even and odd terms
of the Mittag--Leffler series. This check is not an input to any
all-parameter comparison.

\section{Research transition and limits of transfer}
The supplied corpus \cite{Corpus} motivated a search through
positivity, complete monotonicity, and stable-law factorizations.
Its items [116], [122], and [127] concern Le Roy positivity,
stable-law comparisons, and Kanter factors. None of their claimed
theorems is a premise of the proof here.

The useful auxiliary construction problem was to realize all the
original laws on one space, preserve their exact normalization, and
force a positive chance of crossing every fixed threshold. The common
Poisson sequence solves the first task. Norm comparison and the
conditional first-arrival estimate solve the latter tasks. This
construction is applied to the historical comparison itself; an
auxiliary definition is not substituted for the original target.

The positive bound of Theorem~\ref{thm:linear} is a substantive
estimate for the original CDFs. Its proof combines a conditional
inverse threshold, an entropy derivative, Campbell's formula, and
Markov's inequality. These ingredients are standard. Their historical
novelty in this combination remains a separate question.

No estimate for the Riemann hypothesis, the prime-circle global
energy bound, or the Lorch--Muldoon all-orders resultant problem is
proved here. For example, a reciprocal-prime sequence can have
decreasing \(\ell^p\) norms, but it is not a Poisson arrival sequence.
Assigning the stable marginal law to it would require an additional
theorem. Norm monotonicity does not supply the missing arithmetic
cancellation bound.

\appendix
\section{Evidence and the six success conditions}
Mathematical validity and historical priority are separate.
The statements above have complete proofs and were checked
independently within this research session. They have not undergone
external peer review or formal proof-assistant verification.

\small
\begin{longtable}{@{}p{0.065\textwidth}>{\raggedright\arraybackslash}p{0.25\textwidth}>{\raggedright\arraybackslash}p{0.615\textwidth}@{}}
\toprule
Gate & Requirement & Assessment\\
\midrule
\endhead
G1 & Pre-existing identified problem &
Verified from Simon's 2014 Remark~9 and Ferreira--Simon's 2025
Remark~7(b). These author-and-location identifiers are pre-existing;
no universally standardized proper name is assumed.\\[5pt]
G2 & Original statement unchanged &
The exact variables, gamma normalization, probability order, and
full open index interval are retained. Equality in law and the
optimal scale are proved.\\[5pt]
G3 & Actually unresolved &
The cited 2014 and 2025 sources explicitly leave the general
comparison open. This establishes the historical open question
there, not the absence of every intervening resolution worldwide.\\[5pt]
G4 & Unconditional complete proof &
All parameters, thresholds, and scale factors are treated.
No numerical verification or unproved hypothesis is used.
The mathematical argument passed independent internal checks.\\[5pt]
G5 & More than a reformulation &
Norm inequalities, scale optimality, and explicit CDF separation
are proved, including a locally uniform linear bound. The base
principles are classical; a historically new general principle
has not been established.\\[5pt]
G6 & No essentially identical prior solution &
Unverified. Bounded searches and inspection of close precedents
did not locate an identical resolution. Search absence is not
a proof of priority.\\
\bottomrule
\end{longtable}
\normalsize
\noindent\textbf{Overall assessment.}
This is a complete proof draft for the stated stochastic comparison,
with quantitative extensions. It is not a certified all-G1--G6
historic-breakthrough claim.

\subsection{Primary sources and limits of the review}
\begin{enumerate}
\item Simon \cite{Simon}: journal publication and arXiv v2,
especially the definition of optimal stochastic order,
Theorem~3, and Remark~9. Its other stable-power family and
its convex-order comparison are not confused with the target.
\item Ferreira--Simon \cite{FS}: full text, Remark~7(b), and
current arXiv submission history. The latest record inspected
still lists the December 2025 v1.
\item LePage--Woodroofe--Zinn \cite{LWZ}: bibliographic identification
of the classical stable-series lineage. The exact marginal law
used here is derived directly, so no uninspected variant of
their theorem is required.
\item Last--Penrose \cite{LP}: the Laplace functional and Campbell
formula used for the probability estimates.
\item Oesting \cite{Oesting}, Sections~2--3: the prior Poisson
\(\ell^p\) viewpoint and the stable marginal in max-stable
representations. This is a close methodological precedent.
This review does not exhaust all subsequent uses of that viewpoint.
\item Jan\ss en \cite{Janssen}, Section~6 and Corollary~6.3:
stable-power limits of positive sample norms. Combined with norm
monotonicity, such limits offer another route to investigate the
qualitative comparison. The present review therefore makes no claim
that the underlying norm principle is historically new.
\end{enumerate}
Searches included exact article identifiers and combinations of
author names with stochastic order, LePage, Poisson coupling,
Mittag--Leffler, and stable norms. Author publication records
were inspected. The searches establish a bounded review, not
worldwide novelty. The appropriate next independent evaluation
is to check the exact comparison and quantitative estimates
against the stable-law and Poisson-series literature.

\subsection{Correction to the earlier theta-kernel research record}
A preceding note in this session proved infinitely many strict
center failures of odd derivative Laguerre inequalities for the
Riemann Xi kernel while leaving priority unverified. The present
review located Koide's preprint dated 15 August 2026 \cite{Koide}.
Its Theorem~1.1 states a quantitative count already implying
infinitely many such failures. The qualitative conclusion must
therefore be removed from any list of possible first resolutions
produced in this session. This records the existence and scope
of a prior claim; it does not certify every proof in that preprint.

\section{The decisive comparison in compact form}
Take one sequence of rate-one Poisson arrivals and define
\[
X_\alpha=\left(\sum_jT_j^{-1/\alpha}\right)^{-\alpha}.
\]
The strong law makes all sums finite simultaneously for
\(0<\alpha<1\). The Poisson Laplace functional gives
\[
\mathbb E\exp\left(-t\sum_jT_j^{-1/\alpha}\right)
=e^{-\Gamma(1-\alpha)t^\alpha}.
\]
Thus \(X_\alpha\) has the exact law
\(M_\alpha/\Gamma(1-\alpha)\). For every realization,
\(X_\alpha\) is the reciprocal \(\ell^{1/\alpha}\) norm of
the same positive sequence \((1/T_j)\). These reciprocals
strictly decrease with \(\alpha\), proving the original
usual stochastic order. The negative-moment limit proves
optimal scale; Section~\ref{sec:linear} gives the quantitative
bound beyond the basic comparison.

\begin{thebibliography}{9}
\raggedright
\bibitem{Simon}
T.~Simon, \emph{Comparing Fr\'echet and positive stable laws},
Electronic Journal of Probability \textbf{19} (2014), no.~16, 1--25.
\href{https://doi.org/10.1214/EJP.v19-3058}{doi:10.1214/EJP.v19-3058}.
\url{https://arxiv.org/abs/1310.1888}.
\bibitem{FS}
R.~A.~C.~Ferreira and T.~Simon,
\emph{Mittag--Leffler functions and convex ordering},
arXiv:2512.04940v1, 4 December 2025.
\url{https://arxiv.org/abs/2512.04940}.
\bibitem{LWZ}
R.~LePage, M.~Woodroofe, and J.~Zinn,
\emph{Convergence to a stable distribution via order statistics},
Annals of Probability \textbf{9} (1981), no.~4, 624--632.
\href{https://doi.org/10.1214/aop/1176994367}{doi:10.1214/aop/1176994367}.
\bibitem{LP}
G.~Last and M.~Penrose, \emph{Lectures on the Poisson Process},
Cambridge University Press, 2017. Author manuscript:
\url{https://stoch.math.kit.edu/img/Last/lastpenrose2017.pdf}.
\bibitem{Oesting}
M.~Oesting, \emph{Equivalent representations of max-stable
processes via \(\ell^p\) norms}, arXiv:1707.02912, 2017.
\url{https://arxiv.org/abs/1707.02912}.
\bibitem{Janssen}
A.~Jan\ss en, \emph{Limit laws for power sums and norms of i.i.d. samples},
Probability Theory and Related Fields \textbf{146} (2010), 515--533.
\href{https://doi.org/10.1007/s00440-008-0198-y}{doi:10.1007/s00440-008-0198-y}.
\bibitem{Corpus}
QEEDDD, \emph{Verifiable Mathematical Reasoning Corpus},
snapshot inspected 8 November 2026,
commit \texttt{c30d40119134a59eda6fa5cc6331f27c8e67d034}.
\url{https://github.com/QEEDDD/Verifiable-Mathematical-Reasoning-Corpus}.
Research drafts; their claims are not premises of this proof.
\bibitem{Koide}
A.~Koide, \emph{Power-Saving Center-Failure Counts for the
Derivative Laguerre Inequalities of the Riemann Xi Kernel},
preprint, 15 August 2026.
\href{https://doi.org/10.5281/zenodo.21951273}{doi:10.5281/zenodo.21951273}.
\end{thebibliography}
\end{document}
